\section{Aprendizaje de Recompensas con Bradley--Terry y RLHF}

\subsection*{a) Probabilidades, pérdidas y gradientes}

Para un par etiquetado con una respuesta ganadora $y_w$ y una perdedora
$y_l$, el modelo de Bradley--Terry predice la preferencia mediante la
sigmoide de la diferencia de puntajes~\cite{bradley1952,christiano2017}.
Si $z=r_\psi(x,y_w)-r_\psi(x,y_l)$, la probabilidad es $p=\sigma(z)$ y la
pérdida para esa etiqueta es $\mathcal{L}_{\mathrm{RM}}=-\ln p$.

\begin{center}
\begin{tabular}{lrrrr}
    \toprule
    Preferencia & $z$ & $p$ & $\mathcal{L}_{\mathrm{RM}}$ &
    $\lvert \partial\mathcal{L}_{\mathrm{RM}}/\partial z\rvert$ \\
    \midrule
    $A\succ B$ & $2.0-0.5=1.5$ & $0.818$ & $0.201$ & $0.182$ \\
    $A\succ C$ & $2.0-2.5=-0.5$ & $0.378$ & $0.973$ & $0.622$ \\
    \bottomrule
\end{tabular}
\end{center}

La tabla emplea las aproximaciones de probabilidad y logaritmo
indicadas en el enunciado.

Como $\sigma'(z)=\sigma(z)(1-\sigma(z))$, la derivada respecto del
\emph{logit} de preferencia $z$ es
\[
    \frac{\partial\mathcal{L}_{\mathrm{RM}}}{\partial z}
    =-\frac{\sigma'(z)}{\sigma(z)}
    =\sigma(z)-1.
\]
Las derivadas respecto de los puntajes del ganador y del perdedor son
$p-1$ y $1-p$, respectivamente. La comparación $A\succ C$ produce la
mayor magnitud respecto del logit y de los puntajes porque la red predice
que $C$ es más probable como ganador, en desacuerdo con la etiqueta
humana. Para los parámetros de la red se tiene
\[
    \nabla_\psi \mathcal{L}_{\mathrm{RM}}
    =(p-1)\bigl(\nabla_\psi r_\psi(x,y_w)-\nabla_\psi r_\psi(x,y_l)\bigr).
\]
La magnitud de este gradiente también depende de las derivadas de los
puntajes respecto de $\psi$, que el enunciado no especifica.

\subsection*{b) Invariancia ante traslaciones}

Para cualquier $c\in\mathbb{R}$, definimos $r_{\psi'}(x,y)=r_\psi(x,y)+c$.
Al sustituir en la probabilidad,
\begin{align*}
    P_{\psi'}(y_w\succ y_l\mid x)
    &=\sigma\bigl((r_\psi(x,y_w)+c)-(r_\psi(x,y_l)+c)\bigr)\\
    &=\sigma\bigl(r_\psi(x,y_w)-r_\psi(x,y_l)\bigr)
    =P_\psi(y_w\succ y_l\mid x).
\end{align*}
Por tanto, las preferencias identifican diferencias de puntaje y dejan
indeterminado su origen común.

\subsection*{c) Recompensa regularizada de RLHF}

Para $y_1$, el cociente entre la probabilidad de la política actual y
la de la referencia SFT es $0.40/0.10=4$. La recompensa neta es
\[
    R_{\mathrm{RLHF}}(x,y_1)
    =4.0-0.5\ln 4
    \approx 4.0-0.693
    =\boxed{3.307}.
\]
El promedio del logaritmo del cociente bajo $\pi_\theta(\cdot\mid x)$
es $D_{\mathrm{KL}}(\pi_\theta\Vert\pi_{\mathrm{SFT}})$ para ese prompt.
Con $\beta=0$ desaparece el
costo por alejarse de la referencia. La política puede entonces
concentrar probabilidad en respuestas que explotan errores del modelo
de recompensa y reciben puntajes altos pese a ser poco preferidas por
las personas. Este fenómeno se conoce como \emph{reward hacking} o
sobreoptimización de la recompensa aprendida~\cite{ouyang2022}.
